% Bibliografía unificada: todas las entradas y claves de los propietarios.
% Los rótulos de cita se numeran consecutivamente; los textos se preservan.
\providecommand{\etalchar}[1]{$^{#1}$}
\renewcommand{\bibname}{Références bibliographiques}
\begin{thebibliography}{999}
\bibitem{AssmusKey1992}
Edward~F. Assmus, Jr. and Jennifer~D. Key.
\newblock {\em Designs and Their Codes}, volume 103 of {\em Cambridge Tracts in
  Mathematics}.
\newblock Cambridge University Press, 1992.

\bibitem{AbeLamYamada2017}
Toshiyuki Abe, Ching~Hung Lam, and Hiromichi Yamada.
\newblock A remark on {$\mathbb Z_p$}-orbifold constructions of the moonshine
  vertex operator algebra.
\newblock {\em arXiv preprint arXiv:1705.09022}, 2017.
\newblock \href {https://arxiv.org/abs/1705.09022} {\path{arXiv:1705.09022}}.

\bibitem{Apery1979}
Roger Ap{\'e}ry.
\newblock Irrationalit{\'e} de \(\zeta(2)\) et \(\zeta(3)\).
\newblock {\em Ast{\'e}risque}, 61:11--13, 1979.

\bibitem{Apostol1976}
Tom~M. Apostol.
\newblock {\em Introduction to Analytic Number Theory}.
\newblock Springer, 1976.

\bibitem{Araki1976}
H.~Araki.
\newblock Relative entropy of states of von {Neumann} algebras.
\newblock {\em Publications of the Research Institute for Mathematical
  Sciences}, 11:809--833, 1976.

\bibitem{ArchimedesEquilibrium}
Archimedes.
\newblock On the equilibrium of planes.
\newblock The Works of Archimedes, edited by T. L. Heath, 1897.
\newblock URL: \url{https://archive.org/details/worksofarchimede00arch}.

\bibitem{rm:ArakiWoods1968}
H.~Araki and E.~J. Woods.
\newblock A classification of factors.
\newblock {\em Publ. Res. Inst. Math. Sci. Ser. A}, 4:51--130, 1968.
\newblock URL: \url{https://ems.press/content/serial-article-files/41616}.

\bibitem{rm:Bekenstein1973}
J.~D. Bekenstein.
\newblock Black holes and entropy.
\newblock {\em Physical Review D}, 7:2333--2346, 1973.

\bibitem{Bell1964}
John~S. Bell.
\newblock On the einstein podolsky rosen paradox.
\newblock {\em Physics Physique Fizika}, 1(3):195--200, 1964.
\newblock \href {https://doi.org/10.1103/PhysicsPhysiqueFizika.1.195}
  {\path{doi:10.1103/PhysicsPhysiqueFizika.1.195}}.

\bibitem{Barbero}
J.~Fernando Barbero~G.
\newblock Real {Ashtekar} variables for {Lorentzian} signature space-times.
\newblock {\em Physical Review D}, 51:5507--5510, 1995.

\bibitem{Barbero1995}
J.~Fernando Barbero~G.
\newblock Real {Ashtekar} variables for {Lorentzian} signature space-times.
\newblock {\em Physical Review D}, 51:5507--5510, 1995.

\bibitem{Bloch1946}
Felix Bloch.
\newblock Nuclear induction.
\newblock {\em Physical Review}, 70:460--474, 1946.
\newblock \href {https://doi.org/10.1103/PhysRev.70.460}
  {\path{doi:10.1103/PhysRev.70.460}}.

\bibitem{Bochner1955}
Salomon Bochner.
\newblock {\em Harmonic Analysis and the Theory of Probability}.
\newblock University of California Press, Berkeley, 1955.

\bibitem{Born1926}
Max Born.
\newblock Zur quantenmechanik der sto{\ss}vorg{\"a}nge.
\newblock {\em Zeitschrift f{\"u}r Physik}, 37:863--867, 1926.
\newblock \href {https://doi.org/10.1007/BF01397477}
  {\path{doi:10.1007/BF01397477}}.

\bibitem{Borcherds1992}
Richard~E. Borcherds.
\newblock Monstrous moonshine and monstrous lie superalgebras.
\newblock {\em Inventiones Mathematicae}, 109:405--444, 1992.

\bibitem{Bose1924}
Satyendra~Nath Bose.
\newblock Plancks gesetz und lichtquantenhypothese.
\newblock {\em Zeitschrift für Physik}, 26:178--181, 1924.

\bibitem{Bridgman1931}
Percy~W. Bridgman.
\newblock {\em Dimensional Analysis}.
\newblock Yale University Press, 1931.

\bibitem{Buckingham1914}
Edgar Buckingham.
\newblock On physically similar systems; illustrations of the use of
  dimensional equations.
\newblock {\em Physical Review}, 4:345--376, 1914.
\newblock \href {https://doi.org/10.1103/PhysRev.4.345}
  {\path{doi:10.1103/PhysRev.4.345}}.

\bibitem{BIPMSI2019}
{Bureau International des Poids et Mesures}.
\newblock {\em The International System of Units ({SI})}.
\newblock BIPM, 9 edition, 2019.
\newblock Version 3.02.

\bibitem{BIPMSI}
{Bureau International des Poids et Mesures}.
\newblock {\em The International System of Units ({SI} Brochure)}.
\newblock BIPM, S\`evres, 9 edition, 2019.
\newblock Updated online edition.
\newblock \href {https://doi.org/10.59161/AUEZ1291}
  {\path{doi:10.59161/AUEZ1291}}.

\bibitem{birkhoffvonneumann1936}
Garrett Birkhoff and John von Neumann.
\newblock The logic of quantum mechanics.
\newblock {\em Annals of Mathematics}, 37(4):823--843, 1936.

\bibitem{Caratheodory1914}
Constantin Carath{\'e}odory.
\newblock {"U}ber das lineare ma{\ss} von punktmengen---eine verallgemeinerung
  des l{"a}ngenbegriffs.
\newblock {\em Nachrichten von der Gesellschaft der Wissenschaften zu
  G{"o}ttingen, Mathematisch-Physikalische Klasse}, pages 404--426, 1914.

\bibitem{Carnahan2017}
Scott Carnahan.
\newblock 51 constructions of the moonshine module.
\newblock {\em arXiv preprint arXiv:1707.02954}, 2017.
\newblock \href {https://arxiv.org/abs/1707.02954} {\path{arXiv:1707.02954}}.

\bibitem{CGLMP2002}
Daniel Collins, Nicolas Gisin, Noah Linden, Serge Massar, and Sandu Popescu.
\newblock Bell inequalities for arbitrarily high-dimensional systems.
\newblock {\em Physical Review Letters}, 88:040404, 2002.
\newblock \href {https://doi.org/10.1103/PhysRevLett.88.040404}
  {\path{doi:10.1103/PhysRevLett.88.040404}}.

\bibitem{ClauserHorneShimonyHolt1969}
John~F. Clauser, Michael~A. Horne, Abner Shimony, and Richard~A. Holt.
\newblock Proposed experiment to test local hidden-variable theories.
\newblock {\em Physical Review Letters}, 23(15):\allowbreak 880--884, 1969.
\newblock \href {https://doi.org/10.1103/PhysRevLett.23.880}
  {\path{doi:10.1103/PhysRevLett.23.880}}.

\bibitem{CremmerJuliaScherk1978}
Eug\`ene Cremmer, Bernard Julia, and Jo\"el Scherk.
\newblock Supergravity theory in eleven dimensions.
\newblock {\em Physics Letters B}, 76(4):409--412, 1978.
\newblock \href {https://doi.org/10.1016/0370-2693(78)90894-8}
  {\path{doi:10.1016/0370-2693(78)90894-8}}.

\bibitem{ChenLamShimakura2016}
Hsian-Yang Chen, Ching~Hung Lam, and Hiroki Shimakura.
\newblock A {$\mathbb Z_3$}-orbifold construction of the moonshine vertex
  operator algebra and some maximal 3-local subgroups of the monster.
\newblock {\em arXiv preprint arXiv:1606.05961}, 2016.
\newblock \href {https://arxiv.org/abs/1606.05961} {\path{arXiv:1606.05961}}.

\bibitem{ConwayNorton1979}
John~H. Conway and Simon~P. Norton.
\newblock Monstrous moonshine.
\newblock {\em Bulletin of the London Mathematical Society}, 11:308--339, 1979.

\bibitem{CODATA2022}
{CODATA Task Group on Fundamental Constants}.
\newblock The 2022 {CODATA} recommended values of the fundamental physical
  constants, 2024.
\newblock URL: \url{https://physics.nist.gov/constants}.

\bibitem{cohen1963}
Paul~J. Cohen.
\newblock The independence of the continuum hypothesis.
\newblock {\em Proceedings of the National Academy of Sciences of the United
  States of America}, 50(6):1143--1148, 1963.
\newblock \href {https://doi.org/10.1073/pnas.50.6.1143}
  {\path{doi:10.1073/pnas.50.6.1143}}.

\bibitem{cohen1964}
Paul~J. Cohen.
\newblock The independence of the continuum hypothesis, ii.
\newblock {\em Proceedings of the National Academy of Sciences of the United
  States of America}, 51(1):105--110, 1964.
\newblock \href {https://doi.org/10.1073/pnas.51.1.105}
  {\path{doi:10.1073/pnas.51.1.105}}.

\bibitem{cohen1966}
Paul~J. Cohen.
\newblock {\em Set Theory and the Continuum Hypothesis}.
\newblock W. A. Benjamin, New York, 1966.

\bibitem{connes1976}
Alain Connes.
\newblock Classification of injective factors: Cases {II}$_1$, {II}$_\infty$,
  {III}$_\lambda$, $\lambda\neq1$.
\newblock {\em Annals of Mathematics}, 104(1):73--115, 1976.

\bibitem{ConwaySloane1999}
John~H. Conway and Neil J.~A. Sloane.
\newblock {\em Sphere Packings, Lattices and Groups}.
\newblock Springer, 3 edition, 1999.

\bibitem{Dirac1926}
Paul A.~M. Dirac.
\newblock On the theory of quantum mechanics.
\newblock {\em Proceedings of the Royal Society A}, 112:661--677, 1926.

\bibitem{delascuevas2020qms}
Gemma De~las Cuevas, Tom Drescher, and Tim Netzer.
\newblock Quantum magic squares: Dilations and their limitations.
\newblock {\em Journal of Mathematical Physics}, 61(11):111704, 2020.
\newblock \href {https://doi.org/10.1063/5.0022344}
  {\path{doi:10.1063/5.0022344}}.

\bibitem{delascuevas2024simplicial}
Gemma De~las Cuevas, Matt Hoogsteder~Riera, and Tim Netzer.
\newblock Tensor decompositions on simplicial complexes with invariance.
\newblock {\em Journal of Symbolic Computation}, 124:102299, 2024.
\newblock \href {https://doi.org/10.1016/j.jsc.2024.102299}
  {\path{doi:10.1016/j.jsc.2024.102299}}.

\bibitem{delascuevas2021approximate}
Gemma De~las Cuevas, Andreas Klingler, and Tim Netzer.
\newblock Approximate tensor decompositions: disappearance of many separations.
\newblock {\em Journal of Mathematical Physics}, 62:093502, 2021.
\newblock \href {https://doi.org/10.1063/5.0033876}
  {\path{doi:10.1063/5.0033876}}.

\bibitem{delascuevas2023latin}
Gemma De~las Cuevas, Tim Netzer, and Inga Valentiner-Branth.
\newblock Magic squares: Latin, semiclassical and quantum.
\newblock {\em Journal of Mathematical Physics}, 64(2):022201, 2023.
\newblock \href {https://doi.org/10.1063/5.0127393}
  {\path{doi:10.1063/5.0127393}}.

\bibitem{delascuevas2018continuum}
Gemma De~las Cuevas, Norbert Schuch, David P{\'e}rez-Garc{\'i}a, and J.~Ignacio
  Cirac.
\newblock Continuum limits of matrix product states.
\newblock {\em Physical Review B}, 98:174303, 2018.
\newblock \href {https://doi.org/10.1103/PhysRevB.98.174303}
  {\path{doi:10.1103/PhysRevB.98.174303}}.

\bibitem{delescoves2026beyond}
Gemma De~les Coves, Mirte van~der Eyden, and Tim Netzer.
\newblock Beyond operator systems.
\newblock {\em Journal of Mathematical Analysis and Applications}, 556:130102,
  2026.
\newblock \href {https://doi.org/10.1016/j.jmaa.2025.130102}
  {\path{doi:10.1016/j.jmaa.2025.130102}}.

\bibitem{Einstein1918Principles}
Albert Einstein.
\newblock Prinzipielles zur allgemeinen relativitätstheorie.
\newblock {\em Annalen der Physik}, 360(4):241--244, 1918.
\newblock \href {https://doi.org/10.1002/andp.19183600402}
  {\path{doi:10.1002/andp.19183600402}}.

\bibitem{Einstein1924}
Albert Einstein.
\newblock Quantentheorie des einatomigen idealen gases.
\newblock {\em Sitzungsberichte der Preussischen Akademie der Wissenschaften},
  pages 261--267, 1924.

\bibitem{EuclidElements}
Euclid.
\newblock Elements.
\newblock Perseus Digital Library.
\newblock Édition numérique consultée le 24 juillet 2026.
\newblock URL: \url{https://www.perseus.tufts.edu/hopper/text?doc=Euc.+1}.

\bibitem{Everett1957}
III Everett, H.
\newblock Relative state formulation of quantum mechanics.
\newblock {\em Reviews of Modern Physics}, 29:454--462, 1957.

\bibitem{rm:Feigenbaum1978}
Mitchell~J. Feigenbaum.
\newblock Quantitative universality for a class of nonlinear transformations.
\newblock {\em Journal of Statistical Physics}, 19(1):25--52, 1978.
\newblock \href {https://doi.org/10.1007/BF01020332}
  {\path{doi:10.1007/BF01020332}}.

\bibitem{Fermi1926}
Enrico Fermi.
\newblock Sulla quantizzazione del gas perfetto monoatomico.
\newblock {\em Rendiconti Lincei}, 3:145--149, 1926.

\bibitem{Feynman1948}
Richard~P. Feynman.
\newblock Space-time approach to non-relativistic quantum mechanics.
\newblock {\em Reviews of Modern Physics}, 20:367--387, 1948.
\newblock \href {https://doi.org/10.1103/RevModPhys.20.367}
  {\path{doi:10.1103/RevModPhys.20.367}}.

\bibitem{FeynmanHibbs1965}
R.~P. Feynman and A.~R. Hibbs.
\newblock {\em Quantum Mechanics and Path Integrals}.
\newblock McGraw--Hill, New York, 1965.

\bibitem{FLM1988}
Igor Frenkel, James Lepowsky, and Arne Meurman.
\newblock {\em Vertex Operator Algebras and the Monster}.
\newblock Academic Press, 1988.

\bibitem{rm:GibbonsHawking1977}
G.~W. Gibbons and S.~W. Hawking.
\newblock Cosmological event horizons, thermodynamics, and particle creation.
\newblock {\em Physical Review D}, 15:2738--2751, 1977.

\bibitem{gibbons2004phase}
Kathleen~S. Gibbons, Matthew~J. Hoffman, and William~K. Wootters.
\newblock Discrete phase space based on finite fields.
\newblock {\em Physical Review A}, 70:062101, 2004.
\newblock \href {https://doi.org/10.1103/PhysRevA.70.062101}
  {\path{doi:10.1103/PhysRevA.70.062101}}.

\bibitem{Gleason1957}
Andrew~M. Gleason.
\newblock Measures on the closed subspaces of a hilbert space.
\newblock {\em Journal of Mathematics and Mechanics}, 6:885--893, 1957.
\newblock \href {https://doi.org/10.1512/iumj.1957.6.56050}
  {\path{doi:10.1512/iumj.1957.6.56050}}.

\bibitem{glimm1960}
James~G. Glimm.
\newblock On a certain class of operator algebras.
\newblock {\em Transactions of the American Mathematical Society},
  95(2):318--340, 1960.

\bibitem{goedel1931}
Kurt G{\"o}del.
\newblock {\"U}ber formal unentscheidbare s{\"a}tze der principia mathematica
  und verwandter systeme i.
\newblock {\em Monatshefte f{\"u}r Mathematik und Physik}, 38:173--198, 1931.

\bibitem{goedel1940}
Kurt G{\"o}del.
\newblock {\em The Consistency of the Axiom of Choice and of the Generalized
  Continuum-Hypothesis with the Axioms of Set Theory}.
\newblock Number~3 in Annals of Mathematics Studies. Princeton University
  Press, Princeton, NJ, 1940.

\bibitem{Giveon1994}
Amit Giveon, Massimo Porrati, and Eliezer Rabinovici.
\newblock Target space duality in string theory.
\newblock {\em Physics Reports}, 244(2--3):77--202, 1994.
\newblock \href {https://doi.org/10.1016/0370-1573(94)90070-1}
  {\path{doi:10.1016/0370-1573(94)90070-1}}.

\bibitem{gonda2024universality}
Tom{\'a}{\v{s}} Gonda, Tobias Reinhart, Sebastian Stengele, and Gemma De~les
  Coves.
\newblock A framework for universality in physics, computer science, and
  beyond.
\newblock {\em Compositionality}, 6(3), 2024.
\newblock \href {https://doi.org/10.46298/compositionality-6-3}
  {\path{doi:10.46298/compositionality-6-3}}.

\bibitem{rm:Hawking1975}
S.~W. Hawking.
\newblock Particle creation by black holes.
\newblock {\em Communications in Mathematical Physics}, 43:199--220, 1975.

\bibitem{Heisenberg1927}
Werner Heisenberg.
\newblock {\"U}ber den anschaulichen inhalt der quantentheoretischen kinematik
  und mechanik.
\newblock {\em Zeitschrift f{\"u}r Physik}, 43:172--198, 1927.
\newblock \href {https://doi.org/10.1007/BF01397280}
  {\path{doi:10.1007/BF01397280}}.

\bibitem{HornJohnson2013}
Roger~A. Horn and Charles~R. Johnson.
\newblock {\em Matrix Analysis}.
\newblock Cambridge University Press, 2 edition, 2013.

\bibitem{Holst1996}
Sören Holst.
\newblock Barbero's hamiltonian derived from a generalized hilbert--palatini
  action.
\newblock {\em Physical Review D}, 53:5966--5969, 1996.

\bibitem{Hopf1931}
Heinz Hopf.
\newblock {\"U}ber die abbildungen der dreidimensionalen sph{\"a}re auf die
  kugelfl{\"a}che.
\newblock {\em Mathematische Annalen}, 104:637--665, 1931.
\newblock \href {https://doi.org/10.1007/BF01457962}
  {\path{doi:10.1007/BF01457962}}.

\bibitem{HullTownsend1995}
Chris~M. Hull and Paul~K. Townsend.
\newblock Unity of superstring dualities.
\newblock {\em Nuclear Physics B}, 438:109--137, 1995.
\newblock \href {https://arxiv.org/abs/hep-th/9410167}
  {\path{arXiv:hep-th/9410167}}, \href
  {https://doi.org/10.1016/0550-3213(94)00559-W}
  {\path{doi:10.1016/0550-3213(94)00559-W}}.

\bibitem{Hehl1976}
Friedrich~W. Hehl, Paul von~der Heyde, G.~David Kerlick, and James~M. Nester.
\newblock General relativity with spin and torsion: Foundations and prospects.
\newblock {\em Reviews of Modern Physics}, 48:393--416, 1976.

\bibitem{HardyWright2008}
G.~H. Hardy and E.~M. Wright.
\newblock {\em An Introduction to the Theory of Numbers}.
\newblock Oxford University Press, 6 edition, 2008.

\bibitem{IosifescuKraaikamp}
M.~Iosifescu and C.~Kraaikamp.
\newblock {\em Metrical Theory of Continued Fractions}.
\newblock Kluwer, 2002.

\bibitem{Immirzi}
Giorgio Immirzi.
\newblock Real and complex connections for canonical gravity.
\newblock {\em Classical and Quantum Gravity}, 14:L177--L181, 1997.

\bibitem{Immirzi1997}
Giorgio Immirzi.
\newblock Real and complex connections for canonical gravity.
\newblock {\em Classical and Quantum Gravity}, 14:L177--L181, 1997.

\bibitem{Kepler1609}
Johannes Kepler.
\newblock {\em Astronomia Nova}.
\newblock G. Voegelin, Prague, 1609.

\bibitem{Khinchin}
A.~Ya. Khinchin.
\newblock {\em Continued Fractions}.
\newblock University of Chicago Press, 1964.

\bibitem{Kibble1961}
T.~W.~B. Kibble.
\newblock Lorentz invariance and the gravitational field.
\newblock {\em Journal of Mathematical Physics}, 2:212--221, 1961.

\bibitem{Koenig1936}
D{\'e}nes K{"o}nig.
\newblock {\em Theorie der endlichen und unendlichen Graphen}.
\newblock Akademische Verlagsgesellschaft, Leipzig, 1936.

\bibitem{Koblitz1984}
Neal Koblitz.
\newblock {\em p-adic Numbers, p-adic Analysis, and Zeta-Functions}.
\newblock Springer, 2 edition, 1984.

\bibitem{Kolmogorov1933}
Andrei~N. Kolmogorov.
\newblock {\em Grundbegriffe der Wahrscheinlichkeitsrechnung}, volume 2, Heft 3
  of {\em Ergebnisse der Mathematik und ihrer Grenzgebiete}.
\newblock Julius Springer, Berlin, 1933.

\bibitem{Kraus1983}
Karl Kraus.
\newblock {\em States, Effects, and Operations: Fundamental Notions of Quantum
  Theory}, volume 190 of {\em Lecture Notes in Physics}.
\newblock Springer, Berlin, 1983.

\bibitem{Lagarias2013}
Jeffrey~C. Lagarias.
\newblock Euler's constant: Euler's work and modern developments.
\newblock {\em Bulletin of the American Mathematical Society}, 50:527--628,
  2013.

\bibitem{Landauer1961}
Rolf Landauer.
\newblock Irreversibility and heat generation in the computing process.
\newblock {\em IBM Journal of Research and Development}, 5(3):183--191, 1961.
\newblock \href {https://doi.org/10.1147/rd.53.0183}
  {\path{doi:10.1147/rd.53.0183}}.

\bibitem{rm:Lanford1982}
III Lanford, O.~E.
\newblock A computer-assisted proof of the {Feigenbaum} conjectures.
\newblock {\em Bulletin of the American Mathematical Society}, 6:427--434,
  1982.

\bibitem{LindMarcus1995}
Douglas Lind and Brian Marcus.
\newblock {\em An Introduction to Symbolic Dynamics and Coding}.
\newblock Cambridge University Press, 1995.

\bibitem{Lorentz1904}
Hendrik~A. Lorentz.
\newblock Electromagnetic phenomena in a system moving with any velocity
  smaller than that of light.
\newblock {\em Proceedings of the Royal Netherlands Academy of Arts and
  Sciences}, 6:809--831, 1904.

\bibitem{Lothaire2002}
{Lothaire, M.}
\newblock {\em Algebraic Combinatorics on Words}.
\newblock Cambridge University Press, 2002.

\bibitem{Lueders1950}
Gerhart L{\"u}ders.
\newblock {\"U}ber die zustands{\"a}nderung durch den me{\ss}proze{\ss}.
\newblock {\em Annalen der Physik}, 443:322--328, 1950.
\newblock \href {https://doi.org/10.1002/andp.19504430510}
  {\path{doi:10.1002/andp.19504430510}}.

\bibitem{Luders1951}
G.~L{\"u}ders.
\newblock {\"U}ber die zustands{\"a}nderung durch den messprozess.
\newblock {\em Annalen der Physik}, 443:322--328, 1951.

\bibitem{MachMechanics}
Ernst Mach.
\newblock {\em The Science of Mechanics: A Critical and Historical Account of
  Its Development}.
\newblock Open Court, Chicago, 1919.
\newblock English translation of the ninth German edition.

\bibitem{Manivel2005}
Laurent Manivel.
\newblock Configurations of lines and models of lie algebras.
\newblock arXiv:math/0507118, 2005.
\newblock \href {https://arxiv.org/abs/math/0507118}
  {\path{arXiv:math/0507118}}.

\bibitem{martin1975}
Donald~A. Martin.
\newblock Borel determinacy.
\newblock {\em Annals of Mathematics}, 102(2):363--371, 1975.

\bibitem{McKay1980}
John McKay.
\newblock Graphs, singularities, and finite groups.
\newblock In {\em The Santa Cruz Conference on Finite Groups}, volume~37 of
  {\em Proceedings of Symposia in Pure Mathematics}, pages 183--186. American
  Mathematical Society, 1980.

\bibitem{Minkowski1909}
Hermann Minkowski.
\newblock Raum und zeit.
\newblock {\em Physikalische Zeitschrift}, 10:104--111, 1909.

\bibitem{rm:MisnerSharp1964}
C.~W. Misner and D.~H. Sharp.
\newblock Relativistic equations for adiabatic, spherically symmetric
  gravitational collapse.
\newblock {\em Physical Review}, 136:B571--B576, 1964.

\bibitem{MacWilliamsSloane1977}
Florence~Jessie MacWilliams and Neil J.~A. Sloane.
\newblock {\em The Theory of Error-Correcting Codes}.
\newblock North-Holland, 1977.

\bibitem{MaassenUffink1988}
Hans Maassen and J.~B.~M. Uffink.
\newblock Generalized entropic uncertainty relations.
\newblock {\em Physical Review Letters}, 60:1103--1106, 1988.
\newblock \href {https://doi.org/10.1103/PhysRevLett.60.1103}
  {\path{doi:10.1103/PhysRevLett.60.1103}}.

\bibitem{murrayvonneumann1936}
Francis~J. Murray and John von Neumann.
\newblock On rings of operators.
\newblock {\em Annals of Mathematics}, 37(1):116--229, 1936.

\bibitem{NISTTraceability2023}
{National Institute of Standards and Technology}.
\newblock {\em Metrological Traceability, Traceability}.
\newblock National Institute of Standards and Technology, 2023.
\newblock Entrée de glossaire.

\bibitem{DLMFModular}
{National Institute of Standards and Technology}.
\newblock {\em {NIST Digital Library of Mathematical Functions}, Chapter 23:
  Theta Functions and Modular Forms}.
\newblock National Institute of Standards and Technology, 2026.
\newblock Version 1.2.6 ; en particulier les sections 23.15 et 23.18.

\bibitem{Newton1687}
Isaac Newton.
\newblock {\em Philosophiae Naturalis Principia Mathematica}.
\newblock Joseph Streater, London, 1687.

\bibitem{SEPSetTheory}
J.~J. O'Connor and E.~F. Robertson.
\newblock The beginnings of set theory.
\newblock MacTutor History of Mathematics.
\newblock Consulté le 24 juillet 2026.
\newblock URL:
  \url{https://mathshistory.st-andrews.ac.uk/HistTopics/Beginnings_of_set_theory/}.

\bibitem{Parry1964}
W.~Parry.
\newblock Intrinsic markov chains.
\newblock {\em Transactions of the American Mathematical Society}, 112:55--66,
  1964.
\newblock \href {https://doi.org/10.1090/S0002-9947-1964-0161372-1}
  {\path{doi:10.1090/S0002-9947-1964-0161372-1}}.

\bibitem{Pauli1925}
Wolfgang Pauli.
\newblock Über den zusammenhang des abschlusses der elektronengruppen im atom
  mit der komplexstruktur der spektren.
\newblock {\em Zeitschrift für Physik}, 31:765--783, 1925.

\bibitem{planat2007rings}
Michel Planat and Anne-C{\'e}line Baboin.
\newblock Qudits of composite dimension, mutually unbiased bases and projective
  ring geometry.
\newblock {\em Journal of Physics A: Mathematical and Theoretical},
  40:F1005--F1012, 2007.
\newblock \href {https://arxiv.org/abs/0709.2623} {\path{arXiv:0709.2623}}.

\bibitem{Pathria2011}
R.~K. Pathria et Paul~D. Beale.
\newblock {\em Statistical Mechanics}.
\newblock Elsevier, 3 édition, 2011.

\bibitem{rm:Planck}
M.~Planck.
\newblock Ueber das gesetz der energieverteilung im normalspectrum.
\newblock {\em Annalen der Physik}, 309:553--563, 1901.

\bibitem{Planck}
M.~Planck.
\newblock Ueber das gesetz der energieverteilung im normalspectrum.
\newblock {\em Annalen der Physik}, 309:553--563, 1901.

\bibitem{PlanckNobelLecture}
Max Planck.
\newblock The genesis and present state of development of the quantum theory.
\newblock Nobel Lecture, 1920.
\newblock URL:
  \url{https://www.nobelprize.org/prizes/physics/1918/planck/lecture/}.

\bibitem{PoincareScienceHypothesis}
Henri Poincaré.
\newblock {\em La science et l'hypothèse}.
\newblock Flammarion, Paris, 1902.

\bibitem{Robertson1929}
Howard~P. Robertson.
\newblock The uncertainty principle.
\newblock {\em Physical Review}, 34:163--164, 1929.
\newblock \href {https://doi.org/10.1103/PhysRev.34.163}
  {\path{doi:10.1103/PhysRev.34.163}}.

\bibitem{SchrodingerNobel}
Erwin Schrödinger.
\newblock The fundamental idea of wave mechanics.
\newblock Nobel Lecture, 1933.
\newblock URL:
  \url{https://www.nobelprize.org/prizes/physics/1933/schrodinger/lecture/}.

\bibitem{Sciama1964}
Dennis~W. Sciama.
\newblock The physical structure of general relativity.
\newblock {\em Reviews of Modern Physics}, 36:463--469, 1964.

\bibitem{Serre1977Representations}
Jean-Pierre Serre.
\newblock {\em Linear Representations of Finite Groups}.
\newblock Springer, 1977.

\bibitem{Shannon1948}
C.~E. Shannon.
\newblock A mathematical theory of communication.
\newblock {\em Bell System Technical Journal}, 27:379--423, 623--656, 1948.
\newblock \href {https://doi.org/10.1002/j.1538-7305.1948.tb01338.x}
  {\path{doi:10.1002/j.1538-7305.1948.tb01338.x}}.

\bibitem{sion1958}
Maurice Sion.
\newblock On general minimax theorems.
\newblock {\em Pacific Journal of Mathematics}, 8(1):171--176, 1958.

\bibitem{Sommerfeld1916}
Arnold Sommerfeld.
\newblock Zur quantentheorie der spektrallinien.
\newblock {\em Annalen der Physik}, 356(17):1--94, 1916.
\newblock \href {https://doi.org/10.1002/andp.19163561702}
  {\path{doi:10.1002/andp.19163561702}}.

\bibitem{SEPAristotleMath}
{Stanford Encyclopedia of Philosophy}.
\newblock Aristotle and mathematics.
\newblock Consulté le 24 juillet 2026.
\newblock URL: \url{https://plato.stanford.edu/entries/aristotle-mathematics/}.

\bibitem{SEPContinuity}
{Stanford Encyclopedia of Philosophy}.
\newblock Continuity and infinitesimals.
\newblock Consulté le 24 juillet 2026.
\newblock URL: \url{https://plato.stanford.edu/entries/continuity/}.

\bibitem{SEPDedekind}
{Stanford Encyclopedia of Philosophy}.
\newblock Dedekind's contributions to the foundations of mathematics.
\newblock Consulté le 24 juillet 2026.
\newblock URL: \url{https://plato.stanford.edu/entries/dedekind-foundations/}.

\bibitem{Stinespring1955}
William~F. Stinespring.
\newblock Positive functions on c*-algebras.
\newblock {\em Proceedings of the American Mathematical Society}, 6:211--216,
  1955.
\newblock \href {https://doi.org/10.1090/S0002-9939-1955-0069403-4}
  {\path{doi:10.1090/S0002-9939-1955-0069403-4}}.

\bibitem{Stone1932}
Marshall~H. Stone.
\newblock On one-parameter unitary groups in hilbert space.
\newblock {\em Annals of Mathematics}, 33:643--648, 1932.
\newblock \href {https://doi.org/10.2307/1968538} {\path{doi:10.2307/1968538}}.

\bibitem{Susskind1995}
Leonard Susskind.
\newblock The world as a hologram.
\newblock {\em Journal of Mathematical Physics}, 36:6377--6396, 1995.
\newblock \href {https://arxiv.org/abs/hep-th/9409089}
  {\path{arXiv:hep-th/9409089}}.

\bibitem{tHooft1993}
Gerard {'t Hooft}.
\newblock Dimensional reduction in quantum gravity, 1993.
\newblock \href {https://arxiv.org/abs/gr-qc/9310026}
  {\path{arXiv:gr-qc/9310026}}.

\bibitem{PDG2026}
F.~Takahashi et~al.
\newblock Review of particle physics.
\newblock {\em International Journal of Modern Physics A}, 41:2630011, 2026.
\newblock Summary Tables, arrêté éditorial du 15 janvier 2026.

\bibitem{rm:Takesaki}
M.~Takesaki.
\newblock {\em Theory of Operator Algebras II}.
\newblock Springer, 2003.

\bibitem{Tiesinga2021CODATA}
Eite Tiesinga, Peter~J. Mohr, David~B. Newell, and Barry~N. Taylor.
\newblock {CODATA} recommended values of the fundamental physical constants:
  2018.
\newblock {\em Journal of Physical and Chemical Reference Data}, 50(3):033105,
  2021.
\newblock \href {https://doi.org/10.1063/5.0064853}
  {\path{doi:10.1063/5.0064853}}.

\bibitem{Tolar2014}
Ji\v{r}\'i Tolar.
\newblock A classification of finite quantum kinematics, 2014.
\newblock \href {https://arxiv.org/abs/1411.6390} {\path{arXiv:1411.6390}}.

\bibitem{Townsend1995}
Paul~K. Townsend.
\newblock The eleven-dimensional supermembrane revisited.
\newblock {\em Physics Letters B}, 350:184--188, 1995.
\newblock \href {https://arxiv.org/abs/hep-th/9501068}
  {\path{arXiv:hep-th/9501068}}, \href
  {https://doi.org/10.1016/0370-2693(95)00397-4}
  {\path{doi:10.1016/0370-2693(95)00397-4}}.

\bibitem{turing1936}
Alan~M. Turing.
\newblock On computable numbers, with an application to the
  entscheidungsproblem.
\newblock {\em Proceedings of the London Mathematical Society}, 42:230--265,
  1936.

\bibitem{Vardi1988}
Ilan Vardi.
\newblock Determinants of laplacians and multiple gamma functions.
\newblock {\em SIAM Journal on Mathematical Analysis}, 19:493--507, 1988.

\bibitem{vonneumann1923}
John von Neumann.
\newblock Zur einf{\"u}hrung der transfiniten zahlen.
\newblock {\em Acta Litterarum ac Scientiarum Regiae Universitatis Hungaricae
  Francisco-Josephinae, Sectio Scientiarum Mathematicarum}, 1:199--208, 1923.

\bibitem{vonneumann1928}
John von Neumann.
\newblock Zur theorie der gesellschaftsspiele.
\newblock {\em Mathematische Annalen}, 100:295--320, 1928.

\bibitem{vonNeumann1932}
John von Neumann.
\newblock {\em Mathematische Grundlagen der Quantenmechanik}.
\newblock Springer, Berlin, 1932.

\bibitem{vonneumann1932ergodic}
John von Neumann.
\newblock Proof of the quasi-ergodic hypothesis.
\newblock {\em Proceedings of the National Academy of Sciences}, 18(1):70--82,
  1932.

\bibitem{Vourdas2006}
Apostolos Vourdas.
\newblock Quantum systems with finite hilbert space.
\newblock {\em Reports on Progress in Physics}, 67:267--320, 2004.
\newblock Voir également arXiv:quant-ph/0605054 pour les formulations symplectiques de l'espace des phases.

\bibitem{Weyl1927}
H.~Weyl.
\newblock Quantenmechanik und gruppentheorie.
\newblock {\em Zeitschrift f{\"u}r Physik}, 46:1--46, 1927.
\newblock \href {https://doi.org/10.1007/BF02055756}
  {\path{doi:10.1007/BF02055756}}.

\bibitem{rm:Wien}
W.~Wien.
\newblock Eine neue beziehung der strahlung schwarzer k{\"o}rper zum zweiten
  hauptsatz der w{\"a}rmetheorie.
\newblock {\em Sitzungsberichte der K{\"o}niglich Preussischen Akademie der
  Wissenschaften}, pages 55--62, 1893.

\bibitem{Wien}
W.~Wien.
\newblock Eine neue beziehung der strahlung schwarzer k{\"o}rper zum zweiten
  hauptsatz der w{\"a}rmetheorie.
\newblock {\em Sitzungsberichte der K{\"o}niglich Preussischen Akademie der
  Wissenschaften}, pages 55--62, 1893.

\bibitem{Witten1995}
Edward Witten.
\newblock String theory dynamics in various dimensions.
\newblock {\em Nuclear Physics B}, 443:85--126, 1995.
\newblock \href {https://arxiv.org/abs/hep-th/9503124}
  {\path{arXiv:hep-th/9503124}}, \href
  {https://doi.org/10.1016/0550-3213(95)00158-O}
  {\path{doi:10.1016/0550-3213(95)00158-O}}.

\bibitem{HaidaraRamosCierre}
O. Haidara Fall et R. Ramos Balsa,
\emph{El cierre holográfico del infinito: Holografía Modular Triádica y
Mecánica Dimensional}, édition intégrale, 2026.


 Manuscrit inédit ; antécédent documentaire du présent traité. Les constructions et démonstrations utilisées figurent dans le corps de cette édition ; cette entrée en consigne la provenance.
\bibitem{HaidaraRamosMasa}
O. Haidara Fall et R. Ramos Balsa,
\emph{Teoría holográfica integral de la masa HMT--MD}, révision 2, 2026.


 Manuscrit inédit ; antécédent documentaire du présent traité. Les constructions et démonstrations utilisées figurent dans le corps de cette édition ; cette entrée en consigne la provenance.
\bibitem{HaidaraRamosContinuo}
O. Haidara Fall et R. Ramos Balsa,
\emph{La estructura discreta del continuo}, version globale 2, 2026.


 Manuscrit inédit ; antécédent documentaire du présent traité. Les constructions et démonstrations utilisées figurent dans le corps de cette édition ; cette entrée en consigne la provenance.
\bibitem{ArnoldODE}
V. I. Arnold,
\emph{Ordinary Differential Equations},
Springer, troisième édition, 1992.

\bibitem{ArnoldMechanics}
V. I. Arnold,
\emph{Mathematical Methods of Classical Mechanics},
Springer, deuxième édition, 1989.

\bibitem{Cartan}
É. Cartan,
\emph{Leçons sur la géométrie des espaces de Riemann},
Gauthier--Villars, 1928.

\bibitem{DoCarmo}
M. P. do Carmo,
\emph{Differential Geometry of Curves and Surfaces},
Prentice--Hall, 1976.

\bibitem{HilbertCohnVossen}
D. Hilbert et S. Cohn--Vossen,
\emph{Geometry and the Imagination},
Chelsea, 1952.

\bibitem{KatokHasselblatt}
A. Katok et B. Hasselblatt,
\emph{Introduction to the Modern Theory of Dynamical Systems},
Cambridge University Press, 1995.

\bibitem{MilnorDynamics}
J. Milnor,
\emph{Dynamics in One Complex Variable},
Princeton University Press, troisième édition, 2006.

\bibitem{Maxwell}
J. C. Maxwell,
\emph{A Treatise on Electricity and Magnetism},
Clarendon Press, 1873.

\bibitem{Weyl}
H. Weyl,
\emph{The Theory of Groups and Quantum Mechanics},
Dover, 1950.

\bibitem{LandauLifshitzFields}
L. D. Landau et E. M. Lifshitz,
\emph{The Classical Theory of Fields},
Butterworth--Heinemann, quatrième édition, 1975.

\bibitem{x_reciprocidad:feynman1948} R. P. Feynman, ``Space-Time Approach to Non-Relativistic Quantum Mechanics'', \emph{Reviews of Modern Physics} \textbf{20}, 367--387 (1948). \href{https://doi.org/10.1103/RevModPhys.20.367}{doi:10.1103/RevModPhys.20.367}.
\bibitem{x_reciprocidad:feynman1949} R. P. Feynman, ``The Theory of Positrons'', \emph{Physical Review} \textbf{76}, 749--759 (1949). \href{https://doi.org/10.1103/PhysRev.76.749}{doi:10.1103/PhysRev.76.749}.
\bibitem{x_reciprocidad:feynman1965} R. P. Feynman, ``The Development of the Space-Time View of Quantum Electrodynamics'', conférence Nobel, 11 décembre 1965. \href{https://www.nobelprize.org/prizes/physics/1965/feynman/lecture/}{Texte de la conférence}.
\bibitem{x_reciprocidad:mustovicary} B. Musto et J. Vicary, ``Quantum Latin Squares and Unitary Error Bases'', \emph{Quantum Information and Computation} \textbf{16}, 1318--1332 (2016). \href{https://doi.org/10.26421/QIC16.15-16-4}{doi:10.26421/QIC16.15-16-4}.
\bibitem{x_reciprocidad:delascuevas2020} G. De las Cuevas, T. Drescher et T. Netzer, ``Quantum Magic Squares: Dilations and Their Limitations'', \emph{Journal of Mathematical Physics} \textbf{61}, 111704 (2020). \href{https://doi.org/10.1063/5.0022344}{doi:10.1063/5.0022344}.
\bibitem{x_reciprocidad:delascuevas2023} G. De las Cuevas, T. Netzer et I. Valentiner-Branth, ``Magic Squares: Latin, Semiclassical, and Quantum'', \emph{Journal of Mathematical Physics} \textbf{64}, 022201 (2023). \href{https://doi.org/10.1063/5.0127393}{doi:10.1063/5.0127393}.
\bibitem{x_reciprocidad:dirac1933} P. A. M. Dirac, ``The Lagrangian in Quantum Mechanics'', \emph{Physikalische Zeitschrift der Sowjetunion} \textbf{3}, 64--72 (1933). \href{https://www.informationphilosopher.com/solutions/scientists/dirac/Lagrangian_1933.pdf}{Reproduction de l'article original}.
\bibitem{x_reciprocidad:pauli1946} W. Pauli, ``Exclusion Principle and Quantum Mechanics'', conférence Nobel, 13 décembre 1946, dans \emph{Nobel Lectures, Physics 1942--1962}, Elsevier, 1964, pp. 27--43, en particulier p. 42. \href{https://www.nobelprize.org/uploads/2018/06/pauli-lecture.pdf}{Texte de la conférence}.
\bibitem{x_reciprocidad:codata2018} E. Tiesinga, P. J. Mohr, D. B. Newell et B. N. Taylor, ``CODATA Recommended Values of the Fundamental Physical Constants: 2018'', \emph{Reviews of Modern Physics} \textbf{93}, 025010 (2021). \href{https://doi.org/10.1103/RevModPhys.93.025010}{doi:10.1103/RevModPhys.93.025010}. \href{https://physics.nist.gov/cuu/Constants/ArchiveASCII/allascii_2018.txt}{Tableaux archivés de 2018}.
\bibitem{x_reciprocidad:codata2022} P. J. Mohr, D. B. Newell, B. N. Taylor et E. Tiesinga, ``CODATA Recommended Values of the Fundamental Physical Constants: 2022'', \emph{Reviews of Modern Physics} \textbf{97}, 025002 (2025). \href{https://doi.org/10.1103/RevModPhys.97.025002}{doi:10.1103/RevModPhys.97.025002}. \href{https://physics.nist.gov/cuu/Constants/Table/allascii.txt}{Tableau des constantes}.
\bibitem{x_reciprocidad:flm} I. B. Frenkel, J. Lepowsky et A. Meurman, \emph{Vertex Operator Algebras and the Monster}, Pure and Applied Mathematics, vol. 134, Academic Press, 1988.
\bibitem{x_reciprocidad:borcherds} R. E. Borcherds, ``Monstrous Moonshine and Monstrous Lie Superalgebras'', \emph{Inventiones Mathematicae} \textbf{109}, 405--444 (1992). \href{https://doi.org/10.1007/BF01232032}{doi:10.1007/BF01232032}.
\bibitem{x_reciprocidad:conwaysloane} J. H. Conway et N. J. A. Sloane, \emph{Sphere Packings, Lattices and Groups}, 3.\textsuperscript{e} ed., Grundlehren der mathematischen Wissenschaften 290, Springer, 1999. \href{https://doi.org/10.1007/978-1-4757-6568-7}{doi:10.1007/978-1-4757-6568-7}.
\bibitem{x_reciprocidad:bipm2018} Conférence générale des poids et mesures, \emph{Resolution 1 of the 26th CGPM: On the revision of the International System of Units}, 2018. \href{https://www.bipm.org/en/committees/cg/cgpm/26-2018/resolution-1}{Résolution officielle}. \href{https://doi.org/10.59161/CGPM2018RES1E}{doi:10.59161/CGPM2018RES1E}.
\bibitem{x_reciprocidad:dirac1931} P. A. M. Dirac, ``Quantised singularities in the electromagnetic field'', \emph{Proceedings of the Royal Society of London. Series A} \textbf{133}, 60--72 (1931). \href{https://doi.org/10.1098/rspa.1931.0130}{doi:10.1098/rspa.1931.0130}.
\bibitem{x_reciprocidad:lep2006} ALEPH, DELPHI, L3, OPAL et SLD Collaborations; LEP Electroweak Working Group; SLD Electroweak Group et SLD Heavy Flavour Group, ``Precision electroweak measurements on the Z resonance'', \emph{Physics Reports} \textbf{427}, 257--454 (2006). \href{https://doi.org/10.1016/j.physrep.2005.12.006}{doi:10.1016/j.physrep.2005.12.006}.
\bibitem{x_reciprocidad:chenlamshimakura2016} H.-Y. Chen, C. H. Lam et H. Shimakura, ``\(\mathbb Z_3\)-orbifold construction of the Moonshine vertex operator algebra and some maximal \(3\)-local subgroups of the Monster'', prépublication, arXiv:1606.05961 (2016), version 2 (2017). \href{https://arxiv.org/abs/1606.05961}{arXiv:1606.05961}.
\bibitem{x_reciprocidad:abelamyamada2017} T. Abe, C. H. Lam et H. Yamada, ``A remark on \(\mathbb Z_p\)-orbifold constructions of the Moonshine vertex operator algebra'', prépublication, arXiv:1705.09022 (2017), version 4. \href{https://arxiv.org/abs/1705.09022v4}{arXiv:1705.09022}.
\bibitem{x_reciprocidad:kmconwaynorton1979} J. H. Conway et S. P. Norton, ``Monstrous Moonshine'', \emph{Bulletin of the London Mathematical Society} \textbf{11}, 308--339 (1979). \href{https://doi.org/10.1112/blms/11.3.308}{doi:10.1112/blms/11.3.308}.
\bibitem{x_reciprocidad:bennett2003} C. H. Bennett, ``Notes on Landauer's principle, reversible computation, and Maxwell's Demon'', \emph{Studies in History and Philosophy of Modern Physics} \textbf{34}, 501--510 (2003). \href{https://doi.org/10.1016/S1355-2198(03)00039-X}{doi:10.1016/S1355-2198(03)00039-X}.
\bibitem{hmt-bib-nuclear--bibliografia-1}
E. Noether, \emph{Invariante Variationsprobleme}, Nachrichten von der Gesellschaft der Wissenschaften zu Göttingen, Mathematisch-Physikalische Klasse (1918), 235--257. Traduction de M. A. Tavel : \href{https://arxiv.org/abs/physics/0503066}{Invariant Variation Problems
}.
\bibitem{hmt-bib-nuclear--bibliografia-2}
C. Weedbrook, S. Pirandola, R. García-Patrón, N. J. Cerf, T. C. Ralph, J. H. Shapiro et S. Lloyd, \emph{Gaussian Quantum Information}, Reviews of Modern Physics 84 (2012), 621--669. \href{https://arxiv.org/abs/1110.3234}{arXiv:1110.3234}.
\bibitem{hmt-bib-nuclear--bibliografia-3}
J. Williamson, \emph{On the Algebraic Problem Concerning the Normal Forms of Linear Dynamical Systems}, American Journal of Mathematics 58 (1936), 141--163. \href{https://doi.org/10.2307/2371062}{doi:10.2307/2371062}.
\bibitem{hmt-bib-nuclear--bibliografia-4}
A. S. Holevo, M. Sohma et O. Hirota, \emph{Capacity of Quantum Gaussian Channels}, Physical Review A 59 (1999), 1820--1828. \href{https://doi.org/10.1103/PhysRevA.59.1820}{doi:10.1103/PhysRevA.59.1820}.
\bibitem{x_reciprocidad:Manivel2005} L. Manivel, \emph{Configurations of lines and models of Lie algebras}, arXiv:math/0507118, 2005.
\bibitem{x_reciprocidad:Bochner1955} S. Bochner, \emph{Harmonic Analysis and the Theory of Probability}, University of California Press, Berkeley, 1955.
\bibitem{x_reciprocidad:rec:postnikov}
A. Postnikov, \emph{Total positivity, Grassmannians, and networks},
prépublication, 2006. \url{https://arxiv.org/abs/math/0609764}.
\bibitem{ii:bennett2003} C. H. Bennett, ``Notes on Landauer's principle,
reversible computation, and Maxwell's Demon'', \emph{Studies in History
and Philosophy of Modern Physics} \textbf{34}, 501--510 (2003).
\href{https://doi.org/10.1016/S1355-2198(03)00039-X}{doi:10.1016/S1355-2198(03)00039-X}.
\bibitem{ii:feynman1948} R. P. Feynman, ``Space-Time Approach to Non-Relativistic Quantum Mechanics'', \emph{Reviews of Modern Physics} \textbf{20}, 367--387 (1948). \href{https://doi.org/10.1103/RevModPhys.20.367}{doi:10.1103/RevModPhys.20.367}.
\bibitem{ii:mustovicary} B. Musto et J. Vicary, ``Quantum Latin Squares and Unitary Error Bases'', \emph{Quantum Information and Computation} \textbf{16}, 1318--1332 (2016). \href{https://doi.org/10.26421/QIC16.15-16-4}{doi:10.26421/QIC16.15-16-4}.
\bibitem{ii:delascuevas2020} G. De las Cuevas, T. Drescher et T. Netzer, ``Quantum Magic Squares: Dilations and Their Limitations'', \emph{Journal of Mathematical Physics} \textbf{61}, 111704 (2020). \href{https://doi.org/10.1063/5.0022344}{doi:10.1063/5.0022344}.
\bibitem{ii:delascuevas2023} G. De las Cuevas, T. Netzer et I. Valentiner-Branth, ``Magic Squares: Latin, Semiclassical, and Quantum'', \emph{Journal of Mathematical Physics} \textbf{64}, 022201 (2023). \href{https://doi.org/10.1063/5.0127393}{doi:10.1063/5.0127393}.
\bibitem{ii:codata2022} P. J. Mohr, D. B. Newell, B. N. Taylor et E. Tiesinga, ``CODATA Recommended Values of the Fundamental Physical Constants: 2022'', \emph{Reviews of Modern Physics} \textbf{97}, 025002 (2025). \href{https://doi.org/10.1103/RevModPhys.97.025002}{doi:10.1103/RevModPhys.97.025002}. \href{https://physics.nist.gov/cuu/Constants/Table/allascii.txt}{Tableau des constantes}.
\bibitem{ii:flm} I. B. Frenkel, J. Lepowsky et A. Meurman, \emph{Vertex Operator Algebras and the Monster}, Pure and Applied Mathematics, vol. 134, Academic Press, 1988.
\bibitem{ii:borcherds} R. E. Borcherds, ``Monstrous Moonshine and Monstrous Lie Superalgebras'', \emph{Inventiones Mathematicae} \textbf{109}, 405--444 (1992). \href{https://doi.org/10.1007/BF01232032}{doi:10.1007/BF01232032}.
\bibitem{ii:conwaysloane} J. H. Conway et N. J. A. Sloane, \emph{Sphere Packings, Lattices and Groups}, 3.\textsuperscript{e} ed., Grundlehren der mathematischen Wissenschaften 290, Springer, 1999. \href{https://doi.org/10.1007/978-1-4757-6568-7}{doi:10.1007/978-1-4757-6568-7}.
\bibitem{ii:bipm2018} Conférence générale des poids et mesures, \emph{Resolution 1 of the 26th CGPM: On the revision of the International System of Units}, 2018. \href{https://www.bipm.org/en/committees/cg/cgpm/26-2018/resolution-1}{Résolution officielle}. \href{https://doi.org/10.59161/CGPM2018RES1E}{doi:10.59161/CGPM2018RES1E}.
\bibitem{ii:duffokunveneziano} M. J. Duff, L. B. Okun et G. Veneziano, ``Trialogue on the number of fundamental constants'', \emph{JHEP} 03 (2002), 023. \href{https://arxiv.org/abs/physics/0110060}{arXiv:physics/0110060}.
\bibitem{ii:holst1996} S. Holst, ``Barbero's Hamiltonian derived from a generalized Hilbert--Palatini action'', \emph{Physical Review D} \textbf{53} (1996), 5966--5969. \href{https://arxiv.org/abs/gr-qc/9511026}{arXiv:gr-qc/9511026}.
\bibitem{ii:bipmbrochure} Bureau international des poids et mesures, \emph{The International System of Units}, 9.\textsuperscript{e} édition et mises à jour. \href{https://www.bipm.org/en/publications/si-brochure}{Édition officielle}.
\bibitem{iii:feynman1948} R. P. Feynman, ``Space-Time Approach to Non-Relativistic Quantum Mechanics'', \emph{Reviews of Modern Physics} \textbf{20}, 367--387 (1948). \href{https://doi.org/10.1103/RevModPhys.20.367}{doi:10.1103/RevModPhys.20.367}.
\bibitem{iii:feynman1949} R. P. Feynman, ``The Theory of Positrons'', \emph{Physical Review} \textbf{76}, 749--759 (1949). \href{https://doi.org/10.1103/PhysRev.76.749}{doi:10.1103/PhysRev.76.749}.
\bibitem{iii:feynman1965} R. P. Feynman, ``The Development of the Space-Time View of Quantum Electrodynamics'', conférence Nobel, 11 décembre 1965. \href{https://www.nobelprize.org/prizes/physics/1965/feynman/lecture/}{Texte de la conférence}.
\bibitem{iii:mustovicary} B. Musto et J. Vicary, ``Quantum Latin Squares and Unitary Error Bases'', \emph{Quantum Information and Computation} \textbf{16}, 1318--1332 (2016). \href{https://doi.org/10.26421/QIC16.15-16-4}{doi:10.26421/QIC16.15-16-4}.
\bibitem{iii:delascuevas2020} G. De las Cuevas, T. Drescher et T. Netzer, ``Quantum Magic Squares: Dilations and Their Limitations'', \emph{Journal of Mathematical Physics} \textbf{61}, 111704 (2020). \href{https://doi.org/10.1063/5.0022344}{doi:10.1063/5.0022344}.
\bibitem{iii:delascuevas2023} G. De las Cuevas, T. Netzer et I. Valentiner-Branth, ``Magic Squares: Latin, Semiclassical, and Quantum'', \emph{Journal of Mathematical Physics} \textbf{64}, 022201 (2023). \href{https://doi.org/10.1063/5.0127393}{doi:10.1063/5.0127393}.
\bibitem{iii:dirac1933} P. A. M. Dirac, ``The Lagrangian in Quantum Mechanics'', \emph{Physikalische Zeitschrift der Sowjetunion} \textbf{3}, 64--72 (1933). \href{https://www.informationphilosopher.com/solutions/scientists/dirac/Lagrangian_1933.pdf}{Reproduction de l'article original}.
\bibitem{iii:pauli1946} W. Pauli, ``Exclusion Principle and Quantum Mechanics'', conférence Nobel, 13 décembre 1946, dans \emph{Nobel Lectures, Physics 1942--1962}, Elsevier, 1964, pp. 27--43, en particulier p. 42. \href{https://www.nobelprize.org/uploads/2018/06/pauli-lecture.pdf}{Texte de la conférence}.
\bibitem{iii:codata2018} E. Tiesinga, P. J. Mohr, D. B. Newell et B. N. Taylor, ``CODATA Recommended Values of the Fundamental Physical Constants: 2018'', \emph{Reviews of Modern Physics} \textbf{93}, 025010 (2021). \href{https://doi.org/10.1103/RevModPhys.93.025010}{doi:10.1103/RevModPhys.93.025010}. \href{https://physics.nist.gov/cuu/Constants/ArchiveASCII/allascii_2018.txt}{Tableaux archivés de 2018}.
\bibitem{iii:codata2022} P. J. Mohr, D. B. Newell, B. N. Taylor et E. Tiesinga, ``CODATA Recommended Values of the Fundamental Physical Constants: 2022'', \emph{Reviews of Modern Physics} \textbf{97}, 025002 (2025). \href{https://doi.org/10.1103/RevModPhys.97.025002}{doi:10.1103/RevModPhys.97.025002}. \href{https://physics.nist.gov/cuu/Constants/Table/allascii.txt}{Tableau des constantes}.
\bibitem{iii:flm} I. B. Frenkel, J. Lepowsky et A. Meurman, \emph{Vertex Operator Algebras and the Monster}, Pure and Applied Mathematics, vol. 134, Academic Press, 1988.
\bibitem{iii:borcherds} R. E. Borcherds, ``Monstrous Moonshine and Monstrous Lie Superalgebras'', \emph{Inventiones Mathematicae} \textbf{109}, 405--444 (1992). \href{https://doi.org/10.1007/BF01232032}{doi:10.1007/BF01232032}.
\bibitem{iii:conwaysloane} J. H. Conway et N. J. A. Sloane, \emph{Sphere Packings, Lattices and Groups}, 3.\textsuperscript{e} ed., Grundlehren der mathematischen Wissenschaften 290, Springer, 1999. \href{https://doi.org/10.1007/978-1-4757-6568-7}{doi:10.1007/978-1-4757-6568-7}.
\bibitem{iii:bipm2018} Conférence générale des poids et mesures, \emph{Resolution 1 of the 26th CGPM: On the revision of the International System of Units}, 2018. \href{https://www.bipm.org/en/committees/cg/cgpm/26-2018/resolution-1}{Résolution officielle}. \href{https://doi.org/10.59161/CGPM2018RES1E}{doi:10.59161/CGPM2018RES1E}.
\bibitem{iii:dirac1931} P. A. M. Dirac, ``Quantised singularities in the electromagnetic field'', \emph{Proceedings of the Royal Society of London. Series A} \textbf{133}, 60--72 (1931). \href{https://doi.org/10.1098/rspa.1931.0130}{doi:10.1098/rspa.1931.0130}.
\bibitem{iii:lep2006} ALEPH, DELPHI, L3, OPAL et SLD Collaborations; LEP Electroweak Working Group; SLD Electroweak Group et SLD Heavy Flavour Group, ``Precision electroweak measurements on the Z resonance'', \emph{Physics Reports} \textbf{427}, 257--454 (2006). \href{https://doi.org/10.1016/j.physrep.2005.12.006}{doi:10.1016/j.physrep.2005.12.006}.
\bibitem{iv:feynman1948} R. P. Feynman, ``Space-Time Approach to Non-Relativistic Quantum Mechanics'', \emph{Reviews of Modern Physics} \textbf{20}, 367--387 (1948). \href{https://doi.org/10.1103/RevModPhys.20.367}{doi:10.1103/RevModPhys.20.367}.
\bibitem{iv:feynman1949} R. P. Feynman, ``The Theory of Positrons'', \emph{Physical Review} \textbf{76}, 749--759 (1949). \href{https://doi.org/10.1103/PhysRev.76.749}{doi:10.1103/PhysRev.76.749}.
\bibitem{iv:feynman1965} R. P. Feynman, ``The Development of the Space-Time View of Quantum Electrodynamics'', conférence Nobel, 11 décembre 1965. \href{https://www.nobelprize.org/prizes/physics/1965/feynman/lecture/}{Texte de la conférence}.
\bibitem{iv:mustovicary} B. Musto et J. Vicary, ``Quantum Latin Squares and Unitary Error Bases'', \emph{Quantum Information and Computation} \textbf{16}, 1318--1332 (2016). \href{https://doi.org/10.26421/QIC16.15-16-4}{doi:10.26421/QIC16.15-16-4}.
\bibitem{iv:delascuevas2020} G. De las Cuevas, T. Drescher et T. Netzer, ``Quantum Magic Squares: Dilations and Their Limitations'', \emph{Journal of Mathematical Physics} \textbf{61}, 111704 (2020). \href{https://doi.org/10.1063/5.0022344}{doi:10.1063/5.0022344}.
\bibitem{iv:delascuevas2023} G. De las Cuevas, T. Netzer et I. Valentiner-Branth, ``Magic Squares: Latin, Semiclassical, and Quantum'', \emph{Journal of Mathematical Physics} \textbf{64}, 022201 (2023). \href{https://doi.org/10.1063/5.0127393}{doi:10.1063/5.0127393}.
\bibitem{iv:dirac1933} P. A. M. Dirac, ``The Lagrangian in Quantum Mechanics'', \emph{Physikalische Zeitschrift der Sowjetunion} \textbf{3}, 64--72 (1933). \href{https://www.informationphilosopher.com/solutions/scientists/dirac/Lagrangian_1933.pdf}{Reproduction de l'article original}.
\bibitem{iv:pauli1946} W. Pauli, ``Exclusion Principle and Quantum Mechanics'', conférence Nobel, 13 décembre 1946, dans \emph{Nobel Lectures, Physics 1942--1962}, Elsevier, 1964, pp. 27--43, en particulier p. 42. \href{https://www.nobelprize.org/uploads/2018/06/pauli-lecture.pdf}{Texte de la conférence}.
\bibitem{iv:codata2018} E. Tiesinga, P. J. Mohr, D. B. Newell et B. N. Taylor, ``CODATA Recommended Values of the Fundamental Physical Constants: 2018'', \emph{Reviews of Modern Physics} \textbf{93}, 025010 (2021). \href{https://doi.org/10.1103/RevModPhys.93.025010}{doi:10.1103/RevModPhys.93.025010}. \href{https://physics.nist.gov/cuu/Constants/ArchiveASCII/allascii_2018.txt}{Tableaux archivés de 2018}.
\bibitem{iv:codata2022} P. J. Mohr, D. B. Newell, B. N. Taylor et E. Tiesinga, ``CODATA Recommended Values of the Fundamental Physical Constants: 2022'', \emph{Reviews of Modern Physics} \textbf{97}, 025002 (2025). \href{https://doi.org/10.1103/RevModPhys.97.025002}{doi:10.1103/RevModPhys.97.025002}. \href{https://physics.nist.gov/cuu/Constants/Table/allascii.txt}{Tableau des constantes}.
\bibitem{iv:flm} I. B. Frenkel, J. Lepowsky et A. Meurman, \emph{Vertex Operator Algebras and the Monster}, Pure and Applied Mathematics, vol. 134, Academic Press, 1988.
\bibitem{iv:borcherds} R. E. Borcherds, ``Monstrous Moonshine and Monstrous Lie Superalgebras'', \emph{Inventiones Mathematicae} \textbf{109}, 405--444 (1992). \href{https://doi.org/10.1007/BF01232032}{doi:10.1007/BF01232032}.
\bibitem{iv:conwaysloane} J. H. Conway et N. J. A. Sloane, \emph{Sphere Packings, Lattices and Groups}, 3.\textsuperscript{e} ed., Grundlehren der mathematischen Wissenschaften 290, Springer, 1999. \href{https://doi.org/10.1007/978-1-4757-6568-7}{doi:10.1007/978-1-4757-6568-7}.
\bibitem{iv:bipm2018} Conférence générale des poids et mesures, \emph{Resolution 1 of the 26th CGPM: On the revision of the International System of Units}, 2018. \href{https://www.bipm.org/en/committees/cg/cgpm/26-2018/resolution-1}{Résolution officielle}. \href{https://doi.org/10.59161/CGPM2018RES1E}{doi:10.59161/CGPM2018RES1E}.
\bibitem{iv:dirac1931} P. A. M. Dirac, ``Quantised singularities in the electromagnetic field'', \emph{Proceedings of the Royal Society of London. Series A} \textbf{133}, 60--72 (1931). \href{https://doi.org/10.1098/rspa.1931.0130}{doi:10.1098/rspa.1931.0130}.
\bibitem{iv:lep2006} ALEPH, DELPHI, L3, OPAL et SLD Collaborations; LEP Electroweak Working Group; SLD Electroweak Group et SLD Heavy Flavour Group, ``Precision electroweak measurements on the Z resonance'', \emph{Physics Reports} \textbf{427}, 257--454 (2006). \href{https://doi.org/10.1016/j.physrep.2005.12.006}{doi:10.1016/j.physrep.2005.12.006}.
\bibitem{iv:chenlamshimakura2016} H.-Y. Chen, C. H. Lam et H. Shimakura, ``\(\mathbb Z_3\)-orbifold construction of the Moonshine vertex operator algebra and some maximal \(3\)-local subgroups of the Monster'', prépublication, arXiv:1606.05961 (2016), version 2 (2017). \href{https://arxiv.org/abs/1606.05961}{arXiv:1606.05961}.
\bibitem{iv:abelamyamada2017} T. Abe, C. H. Lam et H. Yamada, ``A remark on \(\mathbb Z_p\)-orbifold constructions of the Moonshine vertex operator algebra'', prépublication, arXiv:1705.09022 (2017), version 4. \href{https://arxiv.org/abs/1705.09022v4}{arXiv:1705.09022}.
\bibitem{iv:kmconwaynorton1979} J. H. Conway et S. P. Norton, ``Monstrous Moonshine'', \emph{Bulletin of the London Mathematical Society} \textbf{11}, 308--339 (1979). \href{https://doi.org/10.1112/blms/11.3.308}{doi:10.1112/blms/11.3.308}.
\bibitem{v:mustovicary} B. Musto et J. Vicary, ``Quantum Latin Squares and Unitary Error Bases'', \emph{Quantum Information and Computation} \textbf{16}, 1318--1332 (2016). \href{https://doi.org/10.26421/QIC16.15-16-4}{doi:10.26421/QIC16.15-16-4}.

\bibitem{v:delascuevas2020} G. De las Cuevas, T. Drescher et T. Netzer, ``Quantum Magic Squares: Dilations and Their Limitations'', \emph{Journal of Mathematical Physics} \textbf{61}, 111704 (2020). \href{https://doi.org/10.1063/5.0022344}{doi:10.1063/5.0022344}.

\bibitem{v:delascuevas2023} G. De las Cuevas, T. Netzer et I. Valentiner-Branth, ``Magic Squares: Latin, Semiclassical, and Quantum'', \emph{Journal of Mathematical Physics} \textbf{64}, 022201 (2023). \href{https://doi.org/10.1063/5.0127393}{doi:10.1063/5.0127393}.

\bibitem{v:Bose1924}
Satyendra~Nath Bose.
\newblock Plancks gesetz und lichtquantenhypothese.
\newblock {\em Zeitschrift für Physik}, 26:178--181, 1924.

\bibitem{v:Dirac1926}
Paul A.~M. Dirac.
\newblock On the theory of quantum mechanics.
\newblock {\em Proceedings of the Royal Society A}, 112:661--677, 1926.

\bibitem{v:Einstein1924}
Albert Einstein.
\newblock Quantentheorie des einatomigen idealen gases.
\newblock {\em Sitzungsberichte der Preussischen Akademie der Wissenschaften},
  pages 261--267, 1924.

\bibitem{v:Fermi1926}
Enrico Fermi.
\newblock Sulla quantizzazione del gas perfetto monoatomico.
\newblock {\em Rendiconti Lincei}, 3:145--149, 1926.

\bibitem{v:Pauli1925}
Wolfgang Pauli.
\newblock Über den zusammenhang des abschlusses der elektronengruppen im atom
  mit der komplexstruktur der spektren.
\newblock {\em Zeitschrift für Physik}, 31:765--783, 1925.

\bibitem{v:Pathria2011}
R.~K. Pathria and Paul~D. Beale.
\newblock {\em Statistical Mechanics}.
\newblock Elsevier, 3 edition, 2011.

\bibitem{v:SchrodingerNobel}
Erwin Schrödinger.
\newblock The fundamental idea of wave mechanics.
\newblock Nobel Lecture, 1933.
\newblock URL:
  \url{https://www.nobelprize.org/prizes/physics/1933/schrodinger/lecture/}.

\bibitem{v:Stone1932}
Marshall~H. Stone.
\newblock On one-parameter unitary groups in hilbert space.
\newblock {\em Annals of Mathematics}, 33:643--648, 1932.
\newblock \href {https://doi.org/10.2307/1968538} {\path{doi:10.2307/1968538}}.

\bibitem{v:vonNeumann1932}
John von Neumann.
\newblock {\em Mathematische Grundlagen der Quantenmechanik}.
\newblock Springer, Berlin, 1932.

% BEGIN OWNER /Users/ruben/Documents/ChatGPT/jueces y controles/REVISION_ARGUMENTAL_APERTURAS_CIERRES_20260915/delivery/ARTICLE_V_ES_ARGUMENTAL_20260915/manuscrito/sections/98_referencias_excepcionales_V.tex
\bibitem{v:flm} I. B. Frenkel, J. Lepowsky et A. Meurman, \emph{Vertex Operator Algebras and the Monster}, Pure and Applied Mathematics, vol. 134, Academic Press, 1988.
\bibitem{v:borcherds} R. E. Borcherds, ``Monstrous Moonshine and Monstrous Lie Superalgebras'', \emph{Inventiones Mathematicae} \textbf{109}, 405--444 (1992). \href{https://doi.org/10.1007/BF01232032}{doi:10.1007/BF01232032}.
\bibitem{v:conwaysloane} J. H. Conway et N. J. A. Sloane, \emph{Sphere Packings, Lattices and Groups}, 3.\textsuperscript{e} ed., Grundlehren der mathematischen Wissenschaften 290, Springer, 1999. \href{https://doi.org/10.1007/978-1-4757-6568-7}{doi:10.1007/978-1-4757-6568-7}.


% END OWNER

\bibitem{v:Parry1964}
W.~Parry.
\newblock Intrinsic markov chains.
\newblock {\em Transactions of the American Mathematical Society}, 112:55--66,
  1964.
\newblock \href {https://doi.org/10.1090/S0002-9947-1964-0161372-1}
  {\path{doi:10.1090/S0002-9947-1964-0161372-1}}.
\bibitem{v:bennett2003} C. H. Bennett, ``Notes on Landauer's principle,
reversible computation, and Maxwell's Demon'', \emph{Studies in History
and Philosophy of Modern Physics} \textbf{34}, 501--510 (2003).
\href{https://doi.org/10.1016/S1355-2198(03)00039-X}{doi:10.1016/S1355-2198(03)00039-X}.
\bibitem{vi:mustovicary}
B. Musto et J. Vicary,
``Quantum Latin squares and unitary error bases'',
\emph{Quantum Information and Computation} \textbf{16},
numéros 15--16, 1318--1332 (2016).
\href{https://doi.org/10.26421/QIC16.15-16-4}{doi:10.26421/QIC16.15-16-4}.

\bibitem{vi:delascuevas2020}
G. De las Cuevas, T. Drescher et T. Netzer,
``Quantum magic squares: Dilations and their limitations'',
\emph{Journal of Mathematical Physics} \textbf{61}, 111704 (2020).
\href{https://doi.org/10.1063/5.0022344}{doi:10.1063/5.0022344}.
\href{https://arxiv.org/abs/1912.07332}{Prépublication des auteurs : arXiv:1912.07332}.

\bibitem{vi:delascuevas2023}
G. De las Cuevas, T. Netzer et I. Valentiner-Branth,
``Magic squares: Latin, semiclassical, and quantum'',
\emph{Journal of Mathematical Physics} \textbf{64}, 022201 (2023).
\href{https://doi.org/10.1063/5.0127393}{doi:10.1063/5.0127393}.
\href{https://arxiv.org/abs/2209.10230}{Prépublication des auteurs : arXiv:2209.10230}.

\bibitem{vi:codata2018}
E. Tiesinga, P. J. Mohr, D. B. Newell et B. N. Taylor,
``CODATA recommended values of the fundamental physical constants: 2018'',
\emph{Reviews of Modern Physics} \textbf{93}, 025010 (2021).
\href{https://doi.org/10.1103/RevModPhys.93.025010}{doi:10.1103/RevModPhys.93.025010}.
\href{https://physics.nist.gov/cuu/Constants/archive2018.html}{NIST : archive de l’ajustement de 2018};
\href{https://physics.nist.gov/cuu/Constants/ArchiveASCII/allascii_2018.txt}{tableau ASCII archivé}.

\bibitem{vi:codata2022}
P. J. Mohr, D. B. Newell, B. N. Taylor et E. Tiesinga,
``CODATA recommended values of the fundamental physical constants: 2022'',
\emph{Reviews of Modern Physics} \textbf{97}, 025002 (2025).
\href{https://doi.org/10.1103/RevModPhys.97.025002}{doi:10.1103/RevModPhys.97.025002}.
\href{https://physics.nist.gov/cuu/Constants/article2022.html}{NIST : documentation de l’ajustement de 2022}.
Les comparaisons de l’article conservent identifiés les ajustements de 2018 et 2022.

\bibitem{vi:pdg2026}
F. Takahashi et al. (Particle Data Group),
``Review of Particle Physics'',
\emph{International Journal of Modern Physics A} \textbf{41}, 2630011 (2026).
\href{https://doi.org/10.1142/S0217751X26300115}{doi:10.1142/S0217751X26300115}.
\href{https://pdg.lbl.gov/2026/}{Édition officielle de 2026};
\href{https://pdg.lbl.gov/2026/api/index.html}{documentation d’accès programmatique}.
Le catalogue de cet article conserve l’instantané identifié dans les sources
comme PDG 2026, avec un état documentaire arrêté au 15 janvier 2026 et une consultation consignée
le 26 juillet 2026. Ses 471 enregistrements de masse et 384 de largeur, ainsi que
l’empreinte de provenance, sont maintenus dans le supplément de données ; la référence
web ne remplace pas cet instantané.

\bibitem{viii:mustovicary} B. Musto et J. Vicary, ``Quantum Latin Squares and Unitary Error Bases'', \emph{Quantum Information and Computation} \textbf{16}, 1318--1332 (2016). \href{https://doi.org/10.26421/QIC16.15-16-4}{doi:10.26421/QIC16.15-16-4}.
\bibitem{viii:delascuevas2020} G. De las Cuevas, T. Drescher et T. Netzer, ``Quantum Magic Squares: Dilations and Their Limitations'', \emph{Journal of Mathematical Physics} \textbf{61}, 111704 (2020). \href{https://doi.org/10.1063/5.0022344}{doi:10.1063/5.0022344}.
\bibitem{viii:delascuevas2023} G. De las Cuevas, T. Netzer et I. Valentiner-Branth, ``Magic Squares: Latin, Semiclassical, and Quantum'', \emph{Journal of Mathematical Physics} \textbf{64}, 022201 (2023). \href{https://doi.org/10.1063/5.0127393}{doi:10.1063/5.0127393}.
\bibitem{viii:bombieri2000} E. Bombieri,
``Remarks on Weil's quadratic functional in the theory of prime numbers, I'',
\emph{Rendiconti Lincei. Matematica e Applicazioni},
ser.~9, \textbf{11}, n.~3 (2000), 183--233.
\href{https://www.bdim.eu/item?id=RLIN_2000_9_11_3_183_0}{Biblioteca Digitale Italiana di Matematica}.
\bibitem{ix:feynman1948} R. P. Feynman, ``Space-Time Approach to Non-Relativistic Quantum Mechanics'', \emph{Reviews of Modern Physics} \textbf{20}, 367--387 (1948). \href{https://doi.org/10.1103/RevModPhys.20.367}{doi:10.1103/RevModPhys.20.367}.
\bibitem{ix:mustovicary} B. Musto et J. Vicary, ``Quantum Latin Squares and Unitary Error Bases'', \emph{Quantum Information and Computation} \textbf{16}, 1318--1332 (2016). \href{https://doi.org/10.26421/QIC16.15-16-4}{doi:10.26421/QIC16.15-16-4}.
\bibitem{ix:delascuevas2020} G. De las Cuevas, T. Drescher et T. Netzer, ``Quantum Magic Squares: Dilations and Their Limitations'', \emph{Journal of Mathematical Physics} \textbf{61}, 111704 (2020). \href{https://doi.org/10.1063/5.0022344}{doi:10.1063/5.0022344}.
\bibitem{ix:delascuevas2023} G. De las Cuevas, T. Netzer et I. Valentiner-Branth, ``Magic Squares: Latin, Semiclassical, and Quantum'', \emph{Journal of Mathematical Physics} \textbf{64}, 022201 (2023). \href{https://doi.org/10.1063/5.0127393}{doi:10.1063/5.0127393}.
\bibitem{ix:flm} I. B. Frenkel, J. Lepowsky et A. Meurman, \emph{Vertex Operator Algebras and the Monster}, Pure and Applied Mathematics, vol. 134, Academic Press, 1988.
\bibitem{ix:borcherds} R. E. Borcherds, ``Monstrous Moonshine and Monstrous Lie Superalgebras'', \emph{Inventiones Mathematicae} \textbf{109}, 405--444 (1992). \href{https://doi.org/10.1007/BF01232032}{doi:10.1007/BF01232032}.
\bibitem{ix:conwaysloane} J. H. Conway et N. J. A. Sloane, \emph{Sphere Packings, Lattices and Groups}, 3.\textsuperscript{e} ed., Grundlehren der mathematischen Wissenschaften 290, Springer, 1999. \href{https://doi.org/10.1007/978-1-4757-6568-7}{doi:10.1007/978-1-4757-6568-7}.
\bibitem{ix:chenlamshimakura2016} H.-Y. Chen, C. H. Lam et H. Shimakura, ``\(\mathbb Z_3\)-orbifold construction of the Moonshine vertex operator algebra and some maximal \(3\)-local subgroups of the Monster'', prépublication, arXiv:1606.05961 (2016), version 2 (2017). \href{https://arxiv.org/abs/1606.05961}{arXiv:1606.05961}.
\bibitem{ix:abelamyamada2017} T. Abe, C. H. Lam et H. Yamada, ``A remark on \(\mathbb Z_p\)-orbifold constructions of the Moonshine vertex operator algebra'', prépublication, arXiv:1705.09022 (2017), version 4. \href{https://arxiv.org/abs/1705.09022v4}{arXiv:1705.09022}.
\bibitem{ix:kmconwaynorton1979} J. H. Conway et S. P. Norton, ``Monstrous Moonshine'', \emph{Bulletin of the London Mathematical Society} \textbf{11}, 308--339 (1979). \href{https://doi.org/10.1112/blms/11.3.308}{doi:10.1112/blms/11.3.308}.
\bibitem{ix:Cantor1878} G. Cantor, ``Ein Beitrag zur Mannigfaltigkeitslehre'',
\emph{Journal für die reine und angewandte Mathematik} \textbf{84},
242--258 (1878). \href{https://doi.org/10.1515/crelle-1878-18788413}{doi:10.1515/crelle-1878-18788413}.
\bibitem{ix:Hilbert1902} D. Hilbert, ``Mathematical Problems'',
\emph{Bulletin of the American Mathematical Society} \textbf{8},
437--479 (1902). \href{https://doi.org/10.1090/S0002-9904-1902-00923-3}{doi:10.1090/S0002-9904-1902-00923-3}.
\bibitem{ix:Godel1938} K. Gödel, ``The Consistency of the Axiom of Choice and
of the Generalized Continuum-Hypothesis'', \emph{Proceedings of the National
Academy of Sciences} \textbf{24}, 556--557 (1938).
\href{https://doi.org/10.1073/pnas.24.12.556}{doi:10.1073/pnas.24.12.556}.
\bibitem{ix:Cohen1963} P. J. Cohen, ``The Independence of the Continuum
Hypothesis'', \emph{Proceedings of the National Academy of Sciences}
\textbf{50}, 1143--1148 (1963).
\href{https://doi.org/10.1073/pnas.50.6.1143}{doi:10.1073/pnas.50.6.1143}.
\bibitem{ix:Cohen1964} P. J. Cohen, ``The Independence of the Continuum
Hypothesis, II'', \emph{Proceedings of the National Academy of Sciences}
\textbf{51}, 105--110 (1964).
\href{https://doi.org/10.1073/pnas.51.1.105}{doi:10.1073/pnas.51.1.105}.

\bibitem{ix:Jech2003} T. Jech, \emph{Set Theory}, 3\textsuperscript{e} ed.,
Springer Monographs in Mathematics, Springer, 2003; chapitre 13,
``Constructible Sets'', pp. 175--200.
\href{https://doi.org/10.1007/3-540-44761-X}{doi:10.1007/3-540-44761-X}.

\end{thebibliography}
